\section{Test: Operacions amb polinomis}

\iniciTest{pol1}{b,d,a,c,b,c,d,a,b,a}

\begin{enumerate}

\pregunta Si $p(x)$ i $q(x)$ són dos polinomis, aleshores

\begin{enumerate}[label=\alph*)]

\opcio{a}{$\operatorname{grau}~\left( p(x)+q(x)\right) =$ $\operatorname{grau}~p(x)$ $+~$$\operatorname{grau}~q(x)$}

\opcio{b}{$\operatorname{grau}~\left( p(x)+q(x)\right) \leq \max \left\{ \text{grau}~p(x),\text{grau}~q(x)\right\} $}

\opcio{c}{$\operatorname{grau}~\left( p(x)\cdot q(x)\right) =$ $\operatorname{grau}~p(x)\cdot $$\operatorname{grau}~q(x)$}

\opcio{d}{$\operatorname{grau}~\left( p(x)\cdot q(x)\right) =$ $\operatorname{grau}~\left( p(x)+q(x)\right) $}

\end{enumerate}

\feedback

\pregunta De les expressions següents, quina és un polinomi de $\mathbb{Z}[x]$?

\begin{enumerate}[label=\alph*)]

\opcio{a}{$3x^{4}-2x+\sqrt{2}$}

\opcio{b}{$2x^{-3}+x^{2}-2x+1$}

\opcio{c}{$x^{5}-\frac{8}{2}x^{3}-4x-1$}

\opcio{d}{$3\sqrt{x^{3}}+2x^{2}+x-1$}

\end{enumerate}

\feedback

\pregunta Si $p(x)=2x^{3}-3x^{2}+4x-1$, aleshores

\begin{enumerate}[label=\alph*)]

\opcio{a}{$p(-1)=-4$}

\opcio{b}{$p(\sqrt{2})=-7$}

\opcio{c}{$p(\frac{1}{2})=-\frac{11}{8}$}

\opcio{d}{Cap de les anteriors}

\end{enumerate}

\feedback

\pregunta Sabent que $p(x)=3x^{3}-5x^{2}+5x+1$ i $q(x)=Ax^{3}+(B-C)x^{2}+(2A+B+C)x+1$. Aleshores $p(x)=q(x)$ si

\begin{enumerate}[label=\alph*)]

\opcio{a}{$A=3$, $B=-3$, $C=2$}

\opcio{b}{$A=3$, $B=2$, $C=-3$}

\opcio{c}{$A=3$, $B=C=2$}

\opcio{d}{Cap de les anteriors}

\end{enumerate}

\feedback

\pregunta Donada la potència següent\begin{equation*}
\left( \sqrt{2}+\frac{1}{\sqrt{2}}\right) ^{21} 
\end{equation*}aleshores:

\begin{enumerate}[label=\alph*)]

\opcio{a}{El terme general del seu desenvolupament per la fórmula de Newton és $\binom{21}{k}\left( \sqrt{2}\right) ^{23-2k}$}

\opcio{b}{El terme onzè del seu desenvolupament és $2\binom{21}{10}\sqrt{2}$}

\opcio{c}{L'últim terme del seu desenvolupament és $\frac{1}{1024\sqrt{2}}$}

\opcio{d}{Cap de les anteriors}

\end{enumerate}

\feedback

\pregunta En el desenvolupament per la fórmula de Newton de la potència\begin{equation*}
\left( 2x+\frac{1}{2}\right) ^{15} 
\end{equation*}es compleix:

\begin{enumerate}[label=\alph*)]

\opcio{a}{El terme quart del desenvolupament és $\binom{15}{3}\cdot x^{12}$}

\opcio{b}{El terme independent de $x$ en el desenvolupament és $2^{-15}$}

\opcio{c}{El coeficient del terme de lloc 14 és $\binom{15}{13}\cdot 2^{-9}$}

\opcio{d}{El terme en què apareix $x^{10}$ és $16\cdot \binom{15}{5}\cdot
x^{10}$}

\end{enumerate}

\feedback

\pregunta Calculant els productes indicats 
\begin{equation*}
(x-1+\sqrt{2})\cdot (x-1-\sqrt{2})\cdot (x+1+\sqrt{2})\cdot (x+1-\sqrt{2}) 
\end{equation*}i expressant el resultat en forma de polinomi, s’obté

\begin{enumerate}[label=\alph*)]

\opcio{a}{$x^{4}-6x^{2}-1$}

\opcio{b}{$x^{4}+6x^{2}+1$}

\opcio{c}{$x^{4}-6x^{2}+1$}

\opcio{d}{$x^{4}+6x^{2}-1$}

\end{enumerate}

\feedback

\pregunta El quocient $c(x)$ i el residu $r(x)$ de dividir $12x^{4}-7x^{3}-74x^{2}+7x+10$ per $3x^{2}-7x-4$ són

\begin{enumerate}[label=\alph*)]

\opcio{a}{$c(x)=x^{2}+x+1$ i $r(x)=x-1$}

\opcio{b}{$c(x)=4x^{2}+7x-3$ i $r(x)=7x-1$}

\opcio{c}{$c(x)=x^{2}+x+1$ i $r(x)=2x-2$}

\opcio{d}{$c(x)=4x^{2}+7x-3$ i $r(x)=14x-2$}

\end{enumerate}

\feedback

\pregunta Donats els polinomis\begin{eqnarray*}
p(x) &=&6x^{3}+9x^{2}+2x+3 \\
q(x) &=&2x^{3}+3x^{2}+4x+6
\end{eqnarray*}aleshores:

\begin{enumerate}[label=\alph*)]

\opcio{a}{m.c.d.$\left( p(x),q(x)\right) =x+\frac{3}{2}$}

\opcio{b}{m.c.d.$\left( p(x),q(x)\right) =x^{2}+2$}

\opcio{c}{m.c.d.$\left( p(x),q(x)\right) =x^{3}+\frac{3}{2}x^{2}$}

\opcio{d}{Cap de les anteriors}

\end{enumerate}

\feedback

\pregunta Si $x^{3}-3x^{2}+x-a$ és divisible per $x^{2}+1$, aleshores

\begin{enumerate}[label=\alph*)]

\opcio{a}{$a=4$}

\opcio{b}{$a=3$}

\opcio{c}{$a=2$}

\opcio{d}{$a=1$}

\end{enumerate}

\feedback

\end{enumerate}

\bigskip
